\section{Conclusion}
\label{sec:conclusion}

Four end-to-end driving policies share a failure mode we did not set out to find: when a scene
is re-rendered as night, the representation moves in a direction systematically opposed to the
policy's own hazard-response direction, and the planned speed rises. The effect is carried by
sensor noise rather than by low luminance, and it survives when the perturbation is confined to
the sky---a region that contains nothing a driver needs. Two independent readouts, one
geometric and one behavioural, agree on this in every policy we tested.

We are less able to say what to \emph{do} with the directions themselves. Of five
pre-registered attempts to use their projections as a behavioural proxy---to rank policies, to
flag individual events, to transfer an axis from benchmark to deployment---four are rejected
outright, and the one readout that does transfer has an external target variable rather than a
model-internal one. The direction that is easiest to extract is not the direction that predicts
what the policy will do.

The most severe practical limit is data, not method. Deployment-side hazard events top out in
the low tens under four independent mining criteria, and the public sensor release for the
right-hand-drive city we studied covers $3.9\%$ of its recorded hours. Any claim that a handful
of deployment scenes suffices to estimate a policy's real behaviour has to contend with that
number first.
